\subsection{Forwarding a Reply Without Waiting}

The \code{@ R} field in a script header binds the current reply destination.
It is different from an argument in the message payload.
This forwarding object keeps the original caller's destination:

\sourcelisting[firstline=32]{tutorial/examples/01-basics.abclr2}

\begin{lstlisting}
[forwarder <== [:square 8]]     ; => 64
\end{lstlisting}

The caller waits on the forwarder, but the forwarder's handler does not
wait on the calculator. It sends a past-type message with \code{@ R}.
The calculator's reply then goes directly to the original caller.
Do not add an extra \code{!} to this handler: that would send a second reply.
This style is useful when an object coordinates work but does not need
the result itself.

\subsection{Object Identity and the Current Object}

Inside a script, \code{Me} denotes the current object.
A handler can send a message to itself using \code{[Me <= message]},
usually to arrange work after the current handler finishes.
Under the ordinary metaobject, a handler is processed to completion before
another normal handler for that object begins. This remains true when
the current handler is suspended waiting on a now-type reply.

Consequently, a now-type message to the same ordinary object can deadlock:
the current handler waits for a reply from a handler that cannot begin.
The same problem arises in a cycle: object A waits on B while B waits on A.
Use forwarding or a later past-type phase to break that wait cycle.
Ordinary queue serialization protects the object's handlers from overlapping;
it does not make operations on several objects into one transaction.
A \emph{join} waits for the notifications required by an application
protocol, such as both children's answers, before continuing. It does
not mean that every object or host thread has stopped.

\subsection{Concurrency and Parallel Hardware}

Independent objects can be logically concurrent even when the host executes
one computation at a time. The supplied runtime has a scheduling queue
in \filepath{src/process.lsp}; its existence does not imply that an object
is mapped to a dedicated operating-system thread or CPU.
Long computations in a host Lisp function can also delay other work.
Use the word \emph{concurrency} for the protocol model, and measure the
particular runtime before making claims about parallel speedup.
